\section{从“会操作电脑”到“理解数字任务”}

三位嘉宾的进入路径，恰好说明 GUI Agent 并不是一个突然冒出来的新风口。谷雨从 semantic parsing 讲起：早期研究关注如何把自然语言 instruction 转换成数据库查询、知识库操作或其他 formal representation。那个阶段的环境更可控，任务更像“把一句话翻译成一个可执行结构”。但问题本质已经很接近今天的 computer-using agent：用户给出自然语言目标，系统需要把它转成环境中的可执行动作。

谢天宝的路径则沿着 agent benchmark 与真实 GUI 环境展开。从 AgentBench、intercode、OSWorld 到多模态 GUI grounding，他关注的是怎样让模型在屏幕、网页、操作系统里真正做事，而不是只在文本空间里回答问题。尚晏仪提供了第三条线：从 2021 年开始做能替用户操作电脑的产品，到后来做 Web Agent 和通用 GUI Agent，再到发现通用入口太难商业化，转向 AI 旅行这类具体场景。三条线放在一起看，会发现 GUI Agent 的关键不是“鼠标点击”本身，而是语言、视觉、动作、用户偏好和真实业务流程如何被接到一起。

\begin{knowledgebox}{GUI Agent 的历史前传}
\begin{tabularx}{\textwidth}{>{\raggedright\arraybackslash}p{0.26\textwidth} Y}
\toprule
前传线索 & 对今天 GUI Agent 的意义 \\
\midrule
Semantic Parsing & 训练系统把自然语言任务转成可执行表示，提供了“语言到行动”的早期形式。 \\
Web Automation / RPA & 证明很多数字任务可以被流程化，但也暴露了规则系统对动态环境的脆弱性。 \\
Agent Benchmark & 用 OSWorld、Mind2Web 等任务把能力问题显式化，让模型可以被横向比较。 \\
产品化探索 & 暴露研究 demo 与用户体验之间的距离：能执行不等于好用，能展示过程不等于建立信任。 \\
\bottomrule
\end{tabularx}
\end{knowledgebox}

2025 年的变化，是这些线索终于被更强的多模态模型重新接上。谢天宝提到，早期模型连“点右上角关闭按钮”“把音量拖到 80\%”“修改输入框里的名字”这类看似简单的动作都做不好。后来通过大规模合成数据、环境扩展和模型迭代，GUI grounding 从“纯概念”走向“基本可用”。这不是小修小补，因为 GUI 环境的难点并不只是看见按钮，而是把视觉定位、语言指令、屏幕坐标、动作类型和任务目标同时对齐。

但这也正是产品问题开始变尖锐的地方。尚晏仪讲 Operator 时的态度很典型：它让从业者感到方向被认可，但她并不认为它是一个好产品。用户看着一个小电脑慢慢点击网页，很容易迷茫；一个原本自己几秒钟能完成的动作，被 AI 以很慢的速度展示出来，未必增加信任，反而可能暴露笨拙。研究者想证明“我真的替你操作了”，用户却可能只想知道“结果是否可靠，我是否需要确认”。

\begin{warningbox}{把 GUI Agent 等同于“看见屏幕并点击”，会错过真正的问题}
GUI Agent 的表层动作是点击、输入、滚动和拖拽；但它的深层问题是数字任务理解。一个成熟 Agent 应该知道什么时候点 GUI，什么时候调用 API，什么时候写脚本，什么时候直接读文件，什么时候向用户确认。GUI 是动作空间的一部分，而不是智能本身。
\end{warningbox}

这一章真正建立的是整场对谈的底层前提：GUI Agent 的未来不只是更强 grounding，也不是更逼真的“远程小电脑”。它要解决的是数字环境里的目标理解、行动选择、反馈学习和信任边界。只有把这个前提立住，后面关于 benchmark、商业模式、UDA、memory 和护城河的讨论才会连成一条线。

\subsection{本章小结}

GUI Agent 不是从零出现的新概念，而是 semantic parsing、web automation、agent benchmark 和产品化探索汇合后的结果。2025 年多模态能力让“操作电脑”进入可用区间，但真正难的是从“执行动作”上升到“理解数字任务”。Operator 这类产品的争议说明，研究证明和用户体验之间仍有很长距离。

